% core-motion-sensors.tex — Core Motion: CMMotionManager (one per app), raw vs fused data,
% device axes + attitude reference frames, push vs pull, rates vs battery; the higher-level
% managers (pedometer, activity, altimeter, headphones, batched watch sensors, fall detection),
% authorization, SensorKit in one line.
% Sources (checked 2026-09-29):
%   developer.apple.com/documentation/coremotion/cmmotionmanager (one instance; push/pull; members)
%   developer.apple.com/documentation/coremotion/getting-raw-accelerometer-events (g units, >= 100 Hz)
%   developer.apple.com/documentation/coremotion/getting-raw-gyroscope-events (bias, rad/s)
%   developer.apple.com/documentation/coremotion/getting-processed-device-motion-data (frames, -pi..pi)
%   developer.apple.com/documentation/coremotion/cmattitudereferenceframe · cmdevicemotion/heading
%   developer.apple.com/documentation/coremotion/cmpedometer · cmpedometerdata · cmmotionactivity
%   developer.apple.com/documentation/coremotion/cmaltimeter · cmaltitudedata · cmabsolutealtitudedata
%   developer.apple.com/documentation/coremotion/cmheadphonemotionmanager · cmfalldetectionmanager
%   developer.apple.com/documentation/coremotion/cmbatchedsensormanager + WWDC23 session 10179
%   developer.apple.com/documentation/bundleresources/information-property-list/nsmotionusagedescription
%   developer.apple.com/documentation/sensorkit · developer.apple.com/documentation/updates/coremotion
% @labels: area=ios-platform kind=api platform=ios topic=platform-apis,hardware
% @tags: core-motion, cmmotionmanager, cmdevicemotion, sensor-fusion, attitude, reference-frames, accelerometer, gyroscope, magnetometer, cmpedometer, cmmotionactivitymanager, cmaltimeter
\documentclass[8pt]{extarticle}
\usepackage{printup-sheet}
\usepackage{array}

\lstdefinelanguage{SwiftSheet}{
  morekeywords={class,final,struct,enum,func,var,let,static,private,init,
    if,else,return,guard,self,nil,true,false,in,for},
  sensitive=true, morecomment=[l]{//}, morestring=[b]"}
\lstset{literate={->}{{\hbox{-}\hbox{>}}}2 {==}{{\hbox{=}\hbox{=}}}2}
\lstset{basicstyle=\ttfamily\scriptsize, aboveskip=2pt, belowskip=2pt}
\newcommand\ct[1]{\texttt{#1}}
\newcommand\hd[1]{\par\vspace{3pt}\noindent{\bfseries\color{sheetBlue}#1}\par\vspace{1pt}}

\begin{document}

\sheettitle{Core Motion — sensors, fusion, reference frames}{ios-platform · folio}

\oneliner{Core Motion turns the \textbf{accelerometer} (g, gravity included), \textbf{gyroscope}
(rad/s, biased) and \textbf{magnetometer} ($\mu$T, uncalibrated) into \textbf{device motion} by
\textbf{sensor fusion}: attitude, gravity split from user acceleration, bias-free rotation, calibrated
field. One \ct{CMMotionManager} per app; the motion coprocessor also logs steps and activity
\textbf{while your app is not running} (7 days).}

\vspace{3pt}
\noindent\begin{tikzpicture}[sheet,
    l/.style={font=\tiny, text=black!75, inner sep=1pt, align=center},
    sn/.style={box, font=\scriptsize, text width=27mm, align=left, inner sep=2pt},
    fo/.style={box, font=\scriptsize, text width=51mm, align=left, inner sep=2pt,
                draw=sheetGreen, fill=sheetGreen!8},
    ax/.style={->, very thick, >={Stealth[length=5pt]}}]
  % ---------- the phone and its axes ----------
  \node[font=\bfseries\small, text=sheetBlue, anchor=west] at (-1.75,3.55) {Device axes (portrait)};
  \draw[thick, rounded corners=4pt, fill=black!4] (-0.3,0.5) rectangle (0.9,2.85);
  \draw[sheetGrey, fill=white] (-0.2,0.75) rectangle (0.8,2.6);
  \fill[sheetGrey] (0.3,2.73) circle (0.04);
  \coordinate (o) at (0.3,1.65);
  \draw[ax, sheetRed] (o) -- (1.55,1.65) node[right, font=\scriptsize, text=sheetRed]{\textbf{+x} right};
  \draw[ax, sheetGreen!70!black] (o) -- (0.3,3.3) node[left, font=\scriptsize, text=sheetGreen!60!black]{\textbf{+y} top};
  \draw[ax, sheetBlue] (o) -- (-0.95,0.5) node[below, font=\scriptsize, text=sheetBlue]{\textbf{+z} out of screen};
  \fill (o) circle (0.05);
  \node[l, anchor=north west, align=left] at (1.1,3.1) {\textbf{pitch} = about x\\\textbf{roll} = about y\\\textbf{yaw} = about z};
  \node[l, anchor=north west, align=left, text=sheetBrown] at (1.0,1.35) {flat, face up, still:\\raw accel $\approx$ (0, 0, \textbf{$-$1}) g\\$|a| = 1$ g at rest\\0 g = free fall};
  \node[l, anchor=west, align=left] at (-1.75,-0.25) {rotation: right-hand rule, rad/s\\Euler angles: radians, $-\pi \ldots \pi$};

  \draw[sheetGrey!40] (3.0,3.7) -- (3.0,-0.5);

  % ---------- fusion pipeline ----------
  \node[font=\bfseries\small, text=sheetBlue, anchor=west] at (3.15,3.55) {Raw sensors $\to$ sensor fusion $\to$ \ct{CMDeviceMotion}};
  \node[sn] (acc) at (4.75,2.75) {\textbf{accelerometer}\\\ct{CMAccelerometerData}\\g (9.8 m/s$^2$) · \textbf{gravity + user}};
  \node[sn] (gyr) at (4.75,1.55) {\textbf{gyroscope} \ct{CMGyroData}\\rad/s · \textbf{biased} (temperature)};
  \node[sn] (mag) at (4.75,0.4) {\textbf{magnetometer}\\\ct{CMMagnetometerData}\\$\mu$T · raw, incl. device bias\\{\tiny fused only by the corrected + north frames}};
  \node[box, draw=sheetOrange, fill=sheetOrange!10, font=\scriptsize, text width=16mm, minimum height=28mm] (fus) at (7.85,1.55)
     {\textbf{fusion}\\(motion\\coprocessor +\\Core Motion)\\[3pt]{\tiny frame picked in\\\ct{start\ldots(using:)}}};
  \draw[flow] (acc.east) -- (fus.west |- acc.east);
  \draw[flow] (gyr.east) -- (fus.west |- gyr.east);
  \draw[flow, dashed] (mag.east) -- (fus.west |- mag.east);
  \node[fo] (att) at (11.95,3.0) {\ct{attitude}: \ct{quaternion} · \ct{rotationMatrix} · roll/pitch/yaw\\— relative to the chosen \textbf{reference frame}};
  \node[fo] (grv) at (11.95,2.05) {\ct{gravity} + \ct{userAcceleration} (both in g)\\their sum $=$ raw accel: gravity \textbf{split off} for you};
  \node[fo] (rot) at (11.95,1.1) {\ct{rotationRate} (rad/s) — gyro \textbf{bias removed}};
  \node[fo] (mf) at (11.95,0.25) {\ct{magneticField}: calibrated vector + \ct{accuracy}\\\ct{heading} 0–360° (iOS 11) — \textbf{north frames only}};
  \foreach \t in {att,grv,rot,mf} \draw[hot] (fus.east) -- (\t.west);
  \node[l, anchor=west, align=left, text=sheetRed] at (3.15,-0.35) {\ct{heading} in an arbitrary frame is \textbf{negative} (= invalid), not 0 · \ct{userAcceleration} $\approx$ 0 at rest, whatever the orientation};
\end{tikzpicture}

\begin{multicols}{2}
\raggedright\setstretch{1.0}

\hd{\ct{CMMotionManager}}
\begin{itemize}
  \item \textbf{One instance per app} (Apple: several ``can affect the rate at which the system
        receives data''). Keep a strong reference — a deallocated manager delivers nothing.
  \item Four services, each with \ct{is\ldots Available} / \ct{is\ldots Active}, an
        \ct{\ldots UpdateInterval} in \textbf{seconds} (\ct{1.0/60}), and \ct{start}/\ct{stop}:
        accelerometer, gyro, magnetometer, \textbf{deviceMotion}. Max rate is hardware-dependent,
        ``usually at least \textbf{100 Hz}''; asking for more silently gives the max.
  \item \textbf{Push} = \ct{start\ldots(to: queue) \{ \}}: every sample to an \ct{OperationQueue}.
        \textbf{Pull} = \ct{start\ldots()} with no handler, then read \ct{mm.deviceMotion} once per
        frame (\ct{CADisplayLink}) — Apple's pick for games.
  \item Battery: sensors + fusion run only while started — lowest rate that works, \ct{stop}
        off screen. Updates stop when the app is suspended (no ``motion'' background mode).
  \item \textbf{Prefer deviceMotion} over raw (see picture); \ct{magneticField.accuracy}
        \ct{.uncalibrated}$\ldots$\ct{.high}; \ct{showsDeviceMovementDisplay} lets iOS ask the user to
        move the device to calibrate. \ct{attitude.multiply(byInverseOf: ref)} = rotation
        \emph{since} a reference pose (``recentre'').
\end{itemize}

\hd{Example — fused motion, push, one manager}
\begin{lstlisting}[language=SwiftSheet]
final class Motion {                   // ONE per app, kept alive
  let mm = CMMotionManager()
  private let queue = OperationQueue() // not .main at 100 Hz
  func start(_ onTilt: @escaping (CMAcceleration) -> Void) {
    guard mm.isDeviceMotionAvailable else { return } // Simulator: false
    mm.deviceMotionUpdateInterval = 1.0 / 50         // SECONDS
    let f: CMAttitudeReferenceFrame = .xMagneticNorthZVertical
    guard CMMotionManager.availableAttitudeReferenceFrames()
            .contains(f) else { return }
    mm.startDeviceMotionUpdates(using: f, to: queue) { m, error in
      guard let m else { return }
      onTilt(m.gravity)                // (0,0,-1) when flat
      let jolt = m.userAcceleration    // gravity already removed
      let q = m.attitude.quaternion    // x, y, z, w
      let deg = m.heading              // valid in north frames
    }
  }
  func stop() { mm.stopDeviceMotionUpdates() }
}
\end{lstlisting}

\hd{Authorization}
\ct{NSMotionUsageDescription} is \textbf{required} by \ct{CMPedometer}, \ct{CMMotionActivityManager},
\ct{CMSensorRecorder}, \ct{CMHeadphoneMotionManager} — missing $\to$ \textbf{crash}, not a silent
denial. \ct{CMAuthorizationStatus} (iOS 11): \ct{.notDetermined} \ct{.restricted} \ct{.denied}
\ct{.authorized}. Raw \ct{CMMotionManager}: no status API until \ct{authorizationStatus()} appeared
in the iOS 27.2 \textbf{beta} docs (no discussion yet).

\hd{Attitude reference frames — Z is always vertical}
{\footnotesize
\begin{tabular}{@{}>{\raggedright\arraybackslash}p{33mm}>{\raggedright\arraybackslash}p{42mm}@{}}
\toprule
\textbf{frame} & \textbf{X axis points to} \\ \midrule
\ct{xArbitraryZVertical} & arbitrary — the default \\
\ct{xArbitraryCorrectedZVertical} & arbitrary; magnetometer fixes long-term yaw, more CPU \\
\ct{xMagneticNorthZVertical} & magnetic north — compass \\
\ct{xTrueNorthZVertical} & geographic north — needs location \\
\bottomrule
\end{tabular}}

\columnbreak

\hd{Higher-level managers — use these before raw data}
{\footnotesize
\begin{tabular}{@{}>{\raggedright\arraybackslash}p{21mm}>{\raggedright\arraybackslash}p{54mm}@{}}
\toprule
\ct{CMPedometer} & iOS 8. Steps, distance (m), floors up/down, \ct{currentPace} (\textbf{s/m}),
  \ct{currentCadence} (steps/s). \ct{queryPedometerData(from:to:)} = \textbf{7 days} of history;
  \ct{startUpdates(from:)} live; \ct{is\ldots Available} per metric. \\
\ct{CMMotion\-Activity\-Manager} & iOS 7. \ct{stationary} \ct{walking} \ct{running} \ct{automotive}
  \ct{cycling} \ct{unknown} + \ct{confidence} low/medium/high; live or \ct{queryActivityStarting} (\textbf{7 days}). \\
\ct{CMAltimeter} & iOS 8. Barometer: \ct{relativeAltitude} (m since first event), \ct{pressure}
  (\textbf{kPa}). \textbf{iOS 15}: absolute \ct{altitude} above sea level + \ct{accuracy}. \\
\ct{CMHeadphone\-Motion\-Manager} & iOS 14. Head tracking from spatial-audio AirPods (attitude,
  rotation, user accel); connect/disconnect delegate. \\
\ct{CMBatched\-Sensor\-Manager} & watchOS 10, Series 8 / Ultra, in a \textbf{HealthKit workout}:
  accel \textbf{800 Hz}, device motion \textbf{200 Hz}, one batch per second — golf swings. \\
\ct{CMFall\-Detection\-Manager} & watchOS 7.2, Series 4+: fall events, app woken in background;
  \textbf{Apple-granted entitlement}. \\
\bottomrule
\end{tabular}}\par
\textbf{SensorKit} (iOS 14): research studies only — an entitlement + an approved study; never a
product feature.

\hd{Traps}
\begin{itemize}
  \trap{``Subtract 1 g from z'' works only lying flat — use \ct{userAcceleration} (fused with the gyro).}
  \trap{Interval is \textbf{seconds}, not Hz: \ct{= 60} is one sample a minute.}
  \trap{Handler on \ct{.main} at 100 Hz floods the UI thread — background queue, hop to main.}
  \trap{Euler: \textbf{gimbal lock} at pitch $\pm$90° — use the quaternion.}
  \trap{Steps while the app was closed: no background mode — \emph{query} \ct{CMPedometer} history.}
  \trap{Activity flags are \textbf{not exclusive}: car at a red light = \ct{automotive} +
        \ct{stationary}; results lag ``up to several minutes''.}
  \trap{Shake-to-undo needs no Core Motion: \ct{UIResponder.motionEnded(.motionShake, with:)}.
        Simulator has no motion hardware.}
\end{itemize}

\hd{Remember}
\textbf{``One manager, fused not raw, seconds not Hz, quaternion not Euler — history lives in the
coprocessor.''}


\end{multicols}

\noindent{\footnotesize\color{sheetGrey}\textit{Related:} corelocation-mapkit (heading, true north) ·
background-execution · core-animation-rendering (\ct{CADisplayLink}) · watchos-visionos (workouts) ·
app-hardening-privacy}

\end{document}
